\section{Related Work}
\label{sec:related_work}

Our work draws on several lines of research spanning 3D reconstruction, iterative refinement for geometry estimation, and weight-tied network design.

\parsection{Multi-view 3D reconstruction.}
Classical Structure-from-Motion~\citep{schoenberger2016sfm,pan2024global} recovers geometry through feature matching, pose estimation, and bundle adjustment, but is brittle on in-the-wild scenes with weak texture or dynamic content. Learning-based methods have progressively replaced individual stages of this pipeline, from multi-view stereo~\citep{yao2018mvsnet} to fully end-to-end systems. DUSt3R~\citep{dust3r} and MASt3R~\citep{leroy2025grounding} regress pairwise pointmaps from a CroCo-pretrained~\citep{weinzaepfel2022croco} backbone, while VGGT~\citep{vggt} and concurrent methods~\citep{wang2025pi,lin2025depth,keetha2026mapanything,vggt_omega} process all views jointly through a DINOv2~\citep{oquab2023dinov2}-based transformer, with extensions to incremental capture~\citep{wang2025spann3r,cut3r}, large view counts~\citep{yang2025fast3R, vggttt}, pose-free Gaussian splatting~\citep{ye2025noposplat}, multiple dense geometric quantities~\citep{dens3r} and dynamic scenes~\citep{zhang2025monst3r, vdpm, luo20264rc}. All of these systems use a fixed-depth architecture with many unique parameters. We instead frame multi-view reconstruction as iterative refinement, matching the quality of these deeper feed-forward networks at a fraction of the parameter count while exposing the number of refinement steps as an inference-time compute knob.

\parsection{Iterative refinement.}
RAFT~\citep{RAFT} introduced GRU-based iterative refinement of a per-pixel flow field, an idiom since broadly applied to stereo matching~\citep{lipson2021raft}, visual SLAM~\citep{DroidSLAM,huang2025vipe}, and multi-view stereo~\citep{wang2022itermvs}: in each case, a lightweight recurrent updater iteratively refines geometry on top of a fixed feature backbone. iLRM~\citep{kang2026ilrm} extends iterative refinement to feed-forward 3D Gaussian splatting by treating successive (unshared) transformer layers as optimization steps over a scene representation decoupled from the input views. 
Unlike RAFT-style designs, where a lightweight recurrent updater iterates on a precomputed correspondence volume from a fixed feature backbone, our recurrence interleaves cross-view reasoning and refinement: each step is a full transformer block with frame and global attention, applied to per-view tokens initialized from a DINOv2 patch encoder. Cross-view reasoning and refinement therefore occur within the same repeated computation, rather than being assigned to separate stages.
Unlike iLRM, the block is shared across all $K$ steps and refines the per-view tokens themselves rather than a decoupled scene state.

\parsection{Weight-tied transformers.}
Applying a shared transformer block repeatedly across depth dates back to Universal Transformers~\citep{dehghani2019universal}, which paired weight tying with adaptive halting~\citep{graves2016adaptive}, and ALBERT~\citep{lan2020albert}, which showed that cross-layer parameter sharing yields competitive language models at a fraction of the parameters. Follow-ups have explored weight-tied recurrence in language modeling~\citep{hutchins2022blockrecurrent,yang2024looped,saunshi2025reasoning,geiping2025scaling} and learned-iteration in algorithmic tasks where networks trained for $K$ steps extrapolate to $K' > K$ at test time~\citep{schwarzschild2021algorithm,bansal2022endtoend}.
RAPTOR~\citep{jacobs2025raptor} showed that trained Vision Transformers admit an analogous structure: the layers can be faithfully approximated by a much smaller set of looped blocks, fit via post-hoc distillation against the original network. We adopt RAPTOR's gated block design but apply it differently. Rather than distilling a pretrained network into a looped form, we train a single shared block end-to-end on a 3D reconstruction task loss, with no teacher network and no distillation targets, and surpass an otherwise identical decoupled-parameters variant (\cref{tab:diag_block}) under matched training data and compute.



